\chapter*{前言与致谢}


这是一份面向操作系统课程的教材草稿。它通过研究一个名为 xv6 的示例内核，
来讲解操作系统的主要概念。xv6 以 Dennis Ritchie 和 Ken Thompson 的
Unix Version 6 (v6)~\cite{unix} 为原型。xv6 大体沿用了 v6 的结构和风格，
但使用 ANSI C~\cite{kernighan} 在多核 RISC-V~\cite{riscv} 平台上实现。

建议将本书与 xv6 源代码配合阅读，这种方式受到了 John Lions 的
\textit{Commentary on UNIX 6th Edition}~\cite{lions} 启发；文中给出了
指向源码的超链接：\url{https://github.com/mit-pdos/xv6-riscv}。更多关于
v6 与 xv6 的在线资料（包括若干基于 xv6 的实验作业）可参见
\url{https://pdos.csail.mit.edu/6.1810}。

我们已在 MIT 的操作系统课程 6.828 和 6.1810 中使用这份教材。感谢这些课程
的教师、助教和学生，他们都以直接或间接方式为 xv6 做出了贡献。特别感谢
Adam Belay、Austin Clements 和 Nickolai Zeldovich。最后，感谢所有通过邮件
指出文中错误或提出改进建议的朋友：Abutalib Aghayev, Sebastian Boehm, brandb97,
Anton Burtsev, Raphael Carvalho, Tej Chajed,Brendan Davidson, Rasit Eskicioglu,
Color Fuzzy, Wojciech Gac, Giuseppe, Tao Guo, Haibo Hao, Naoki Hayama, Chris
Henderson, Robert Hilderman, Eden Hochbaum, Wolfgang Keller, Paweł Kraszewski,
Henry Laih, Jin Li, Austin Liew, lyazj@github.com, Pavan Maddamsetti, Jacek
Masiulaniec, Michael McConville, m3hm00d, Mes0903, miguelgvieira, Mark Morrissey,
Muhammed Mourad, Harry Pan, Harry Porter, Siyuan Qian, Zhefeng Qiao, Askar Safin,
Salman Shah, Huang Sha, Vikram Shenoy, Adeodato Simó, Ruslan Savchenko, Pawel
Szczurko, Warren Toomey, tyfkda, tzerbib, Vanush Vaswani, Chen Wang, Xi Wang, and
Zou Chang Wei, Sam Whitlock, Qiongsi Wu, LucyShawYang, ykf1114@gmail.com, and
Meng Zhou

如果你发现错误或有改进建议，请邮件联系 Frans Kaashoek 和 Robert Morris
（kaashoek,rtm@csail.mit.edu）。
